% !TEX program = pdflatex
\documentclass[12pt,a4paper]{article}
\usepackage[T1]{fontenc}
\usepackage{lmodern}
\usepackage{amsmath,amssymb}
\usepackage{geometry}
\geometry{margin=2.5cm}
\usepackage{hyperref}
\hypersetup{colorlinks=true,linkcolor=blue,citecolor=blue,urlcolor=blue}
\usepackage{booktabs}

\title{Information Transmission Model via Delay Circuits:\\Summary of 5 Papers}
\author{Noriaki Kihara\\
WF System Co., Ltd.\\
Osaka University, Faculty of Engineering Science (graduated)\\
ORCID: 0009-0004-6753-4020}
\date{April 2026}

\begin{document}
\maketitle

\section*{Paper List}

\begin{table}[h]
\centering
\begin{tabular}{clc}
\toprule
\# & Title & DOI \\
\midrule
{[1]} & An Information-Theoretic Framework for Information Transmission & \href{https://doi.org/10.5281/zenodo.19534345}{10.5281/zenodo.19534345} \\
{[2]} & Symmetries Inherent in the Delay Circuit Model & \href{https://doi.org/10.5281/zenodo.19534347}{10.5281/zenodo.19534347} \\
{[3]} & Realization of Simple Harmonic Oscillation and Sinusoidal Waves & \href{https://doi.org/10.5281/zenodo.19534349}{10.5281/zenodo.19534349} \\
{[4]} & Realization of Perfectly Elastic Collision & \href{https://doi.org/10.5281/zenodo.19534353}{10.5281/zenodo.19534353} \\
{[5]} & Realization of a Wave Packet Collapse Model & \href{https://doi.org/10.5281/zenodo.19534355}{10.5281/zenodo.19534355} \\
\bottomrule
\end{tabular}
\end{table}

\section*{Purpose of This Research}

Starting from only two computational elements --- the delay circuit $F_D$ (which transmits information with a one-step delay) and the NOT operation circuit (which inverts information instantaneously) --- we deductively construct patterns ranging from oscillation, propagation, and collision to wave packet collapse. The starting point is the minimal definition that ``the transmission of information $I$ is accompanied by a delay.'' No physical claims are made.

\section*{Constructed Computational Propositions}

\subsection*{Paper [1]: Basic Structure of Information Transmission}

\textbf{1. Three Properties of Information $I$}\quad
Three properties are defined for information $I$: existence of a null state, equivalence judgment, and inversion ($\overline{\overline{I}} = I$). (\S 2)

\textbf{2. Delay Circuit $F_D$ and NOT Operation Circuit}\quad
$F_D$ is defined as a circuit that outputs input information after a one-step delay. The NOT operation circuit is defined as a circuit that inverts input information without delay. (\S 3)

\textbf{3. Simple Harmonic Oscillation Pattern}\quad
By feedback connection of $F_D$ and NOT, the oscillation pattern $x(n) = I \cdot (-1)^{n+1}$ is constructed. (\S 4.1)

\textbf{4. Constant-Velocity Propagation Pattern}\quad
By serial connection of $F_D$ (open string model), the linear propagation pattern with arrival position $k = n$ is constructed. (\S 4.2)

\textbf{5. Fundamental String Vibration Pattern}\quad
By the closed string model ($F_D$ serial + NOT), a fundamental vibration pattern where one oscillation $= 2n$ operations is constructed. (\S 4.3)

\subsection*{Paper [2]: Framework Extension}

\textbf{6. Inversion Symmetry}\quad
The operation rule of $F_D$ is invariant under the transformation $I \to \overline{I}$. (\S 3.1)

\textbf{7. Superclass and Meta-Information}\quad
A superclass for grouping information of different classes, and meta-information (counter $n$, operation function $*$) attached to information are defined. (\S 3.2, \S 3.3)

\textbf{8. External Intervention}\quad
An operation that forcibly rewrites the internal information of a specific $F_D$ from outside the system is defined. From the perspective of the system's interior, this is observed as an event that suddenly appears from outside causality. (\S 3.5)

\subsection*{Paper [3]: Four Types of Wave Patterns}

\textbf{9. Class $A$ (Square Wave)}\quad
A square wave is constructed as a class without an operation function $*$. (\S 3.3)

\textbf{10. Class $B$ (Transverse Wave / Sinusoidal Wave)}\quad
By setting $w = w_{\max} \cdot \sin(n/m \cdot 360^\circ)$ as the operation function $*$, a transverse wave with displacement perpendicular to the propagation direction is constructed. (\S 3.4)

\textbf{11. Class $C$ (Longitudinal Wave)}\quad
Using the same $*$ as Class $B$, a longitudinal wave is constructed by interpreting the output $w$ as displacement in the propagation direction. (\S 3.5)

\textbf{12. Class $D$ (Rotational Wave)}\quad
By setting $w = (n/m \cdot 360^\circ) \bmod 360^\circ$ as $*$, a rotational wave is constructed. (\S 3.6)

The definitions of $F_D$ and the NOT operation circuit are not modified at all. Differences in wave types are expressed solely through the choice of class definition and $*$.

\subsection*{Paper [4]: Perfectly Elastic Collision}

\textbf{13. Class $E$ (Transverse Wave with Position Information)}\quad
Class $B$ is extended to define Class $E$ with initial position $x_0$, current position $x$, and displacement $dx$. (\S 3)

\textbf{14. Non-Interaction Model (Pass-Through)}\quad
A model is constructed where two instances of Class $E$ propagate in opposite directions and pass through each other without interaction at the same position. (\S 4)

\textbf{15. Class $F$ (with Collision Functionality)}\quad
Class $F$ is defined by adding collision radius $r$ and collision condition $(x_1 - x_2)^2 < \max(r_1^2, r_2^2)$ to Class $E$. Using $w_{\max}$ as an analogue of mass, $dx$ is updated according to the formula for one-dimensional perfectly elastic collision. (\S 5)

\textbf{16. Conservation of Momentum and Energy}\quad
It is verified that the momentum analogue $Q = w_{\max,1} \cdot dx_1 + w_{\max,2} \cdot dx_2$ and the energy analogue $K = w_{\max,1} \cdot dx_1^2 + w_{\max,2} \cdot dx_2^2$ are conserved before and after collision. (\S 5.2)

\subsection*{Paper [5]: Wave Packet Collapse}

\textbf{17. Class $G$ (Promotion of Wave Packet Radius and Maximum Wave Height to Information)}\quad
Class $G$ is defined by promoting the meta-information $r$ (wave packet radius) and $w_{\max}$ (maximum wave height) of Class $F$ to information. (\S 2)

\textbf{18. Mass Analogue $M = w_{\max} \cdot r$}\quad
The mass analogue is defined as $M = w_{\max} \cdot r$, and it is shown that $M$ is conserved before and after collision. (\S 2.2)

\textbf{19. Wave Packet Collapse Rule}\quad
A rule is constructed where the wave packet radius collapses as $r^{\text{new}} = \min(r_1, r_2)$ upon collision, and the maximum wave height increases as $w_{\max}^{\text{new}} = M / r^{\text{new}}$ due to conservation of $M$. (\S 2.4)

\textbf{20. Wave Packet Collapse Sequence}\quad
A process is realized where a photon with a wide wave packet ($r = 48$, $w_{\max} = 1$) collides with an electron with a narrow wave packet ($r = 6$), causing the wave packet to collapse to $r = 6$, $w_{\max} = 8$. (\S 3)

\section*{What This Research Does Not Claim}

This research is limited to describing the consistency of the computational model. The following claims are not made:

\begin{itemize}
  \item That delay circuits represent physical entities or devices
  \item That the derived patterns constitute derivations or proofs of physical laws
  \item That the physical mechanism of wave packet collapse is explained
  \item That the mass analogue $w_{\max}$ is identical to physical mass
  \item That external intervention corresponds to quantum mechanical observation
\end{itemize}

These physical identifications require additional assumptions beyond the computational propositions of this research (Paper [5] \S 4).

\section*{References}

{[1]} Kihara, N. (2026). \textit{An Information-Theoretic Framework for Information Transmission}. DOI: \href{https://doi.org/10.5281/zenodo.19534345}{10.5281/zenodo.19534345}.

{[2]} Kihara, N. (2026). \textit{Symmetries Inherent in the Delay Circuit Model}. DOI: \href{https://doi.org/10.5281/zenodo.19534347}{10.5281/zenodo.19534347}.

{[3]} Kihara, N. (2026). \textit{Realization of Simple Harmonic Oscillation and Sinusoidal Waves in the Delay Circuit Model}. DOI: \href{https://doi.org/10.5281/zenodo.19534349}{10.5281/zenodo.19534349}.

{[4]} Kihara, N. (2026). \textit{Realization of Perfectly Elastic Collision in the Delay Circuit Model}. DOI: \href{https://doi.org/10.5281/zenodo.19534353}{10.5281/zenodo.19534353}.

{[5]} Kihara, N. (2026). \textit{Realization of a Wave Packet Collapse Model in the Delay Circuit Model}. DOI: \href{https://doi.org/10.5281/zenodo.19534355}{10.5281/zenodo.19534355}.

\end{document}
